\clearpage\SourceLeaf{462}
\begin{center}\Large\textls[200]{INHALT.}\par\medskip\large GAUSS WERKE BAND VIII. NACHTRÄGE ZU BAND I—IV.\end{center}
\IndexHeading{ARITHMETIK UND ALGEBRA.}
\IndexSubheading{Nachlass.}
\IndexRow{0}{Zwei Notizen über die Auflösung der Congruenz $xx+yy+zz\equiv0\ (\mathrm{mod.}\,p)$}{.\,.\,.\,.\,.}{Seite}{3}
\IndexRow{0}{Notizen über cubische und biquadratische Reste}{.\,.\,.\,.\,.\,.\,.\,.\,.\,.\,.\,.\,.\,.\,.\,.}{—}{5}
\IndexRow{0}{Zur Theorie der cubischen Reste}{.\,.\,.\,.\,.\,.\,.\,.\,.\,.\,.\,.\,.\,.\,.\,.\,.\,.\,.\,.\,.}{—}{15}
\IndexRow{0}{Fragmente zur Theorie der aus einer Cubikwurzel zu bildenden ganzen algebraischen Zahlen}{}{—}{21}
\IndexRow{0}{Beweis der Irrationalität der Tangenten rationaler Bögen in einer neuen Gestalt}{.\,.\,.\,.}{—}{27}
\IndexRow{0}{Notiz über Auflösung eines speciellen Systems linearer Gleichungen}{.\,.\,.\,.\,.\,.\,.\,.\,.}{—}{30}
\IndexRow{0}{Mechanischer Satz über die Wurzeln einer ganzen Function $f(x)$ und ihrer Ableitung $f^{\prime}(x)$}{}{—}{32}
\IndexHeading{ANALYSIS UND FUNCTIONENTHEORIE.}
\IndexSubheading{Nachlass.}
\IndexRow{0}{De integratione formulae differentialis $(1+n\cos\varphi)^{\nu}d\varphi$}{.\,.\,.\,.\,.\,.\,.\,.\,.\,.\,.\,.\,.\,.}{—}{35}
\IndexRow{0}{Beweis eines von \textsc{Euler} aufgestellten Satzes über exacte Differentialausdrücke}{.\,.\,.\,.\,.}{—}{65}
\IndexRow{0}{Vier Notizen über Inversion der Potenzreihen}{.\,.\,.\,.\,.\,.\,.\,.\,.\,.\,.\,.\,.\,.\,.\,.\,.}{—}{69}
\IndexRow{0}{Neuer Beweis des \textsc{Lagrange}schen Lehrsatzes}{.\,.\,.\,.\,.\,.\,.\,.\,.\,.\,.\,.\,.\,.\,.\,.\,.}{—}{76}
\IndexRow{0}{\textsc{Lagranges} Lehrsatz, auf möglich lichtvollste Art abgeleitet}{.\,.\,.\,.\,.\,.\,.\,.\,.\,.\,.}{—}{80}
\IndexRow{0}{Entwickelung von $\dfrac{1}{(h-\cos\varphi)^{\frac32}}$ in die Reihe $A^{(0)}+2A^{(1)}\cos\varphi+2A^{(2)}\cos2\varphi+\text{etc.}$}{.\,.\,.\,.}{—}{84}
\IndexRow{0}{Schönes Theorem der Wahrscheinlichkeitsrechnung}{.\,.\,.\,.\,.\,.\,.\,.\,.\,.\,.\,.\,.\,.\,.}{—}{88}
\IndexRow{0}{Über das Wesen und die Definition der Functionen}{.\,.\,.\,.\,.\,.\,.\,.\,.\,.\,.\,.\,.\,.}{—}{90}
\IndexRow{0}{Untersuchungen über die transcendenten Functionen, die aus dem Integral $\displaystyle\int\dfrac{dx}{\sqrt{1+x^3}}$ ihren Ursprung haben}{.\,.\,.\,.\,.\,.\,.\,.\,.\,.\,.\,.\,.\,.\,.\,.\,.\,.\,.\,.\,.\,.\,.\,.\,.\,.}{—}{93}
\IndexRow{0}{Inversion des elliptischen Integrals erster Gattung}{.\,.\,.\,.\,.\,.\,.\,.\,.\,.\,.\,.\,.\,.\,.}{—}{96}
\IndexRow{0}{Theorema elegantissimum [das arithmetisch-geometrische Mittel betreffend]}{.\,.\,.\,.\,.\,.}{—}{98}
\IndexRow{0}{Drei Fragmente über elliptische Modulfunctionen}{.\,.\,.\,.\,.\,.\,.\,.\,.\,.\,.\,.\,.\,.\,.}{—}{99}
\IndexRow{0}{Weitere Fragmente über das Pentagramma mirificum}{.\,.\,.\,.\,.\,.\,.\,.\,.\,.\,.\,.\,.\,.}{—}{106}
